% ===================================================================
%  Dodatek V: Formalizacja sektora SU(3)_c — emergencja gluonów,
%             konfinowanie i pełna struktura cechowania SM
%  TGP v1 — Sesja v42 (2026-04-01)
% ===================================================================
\section{Sektor $\mathrm{SU}(3)_c$: emergencja gluon\'ow
  i~konfinowanie kolor\'ow}%
\label{app:V-su3}

Niniejszy dodatek formalizuje wyniki sek.~\ref{ssec:su3-color}
dotycz\k{a}ce sektora kolorowego TGP.
W~po\l\k{a}czeniu z~dod.~\ref{app:u1_formal}
(sektor $\mathrm{U}(1)$) i~dod.~\ref{app:U-su2} (sektor
$\mathrm{SU}(2)\times\mathrm{U}(1)$), niniejszy dodatek
domyka \textbf{pe\l{}n\k{a} struktur\k{e} cechowania TGP}:
\begin{equation}\label{eq:V-gauge-group}
  G_{\rm SM} \;=\;
  \underbrace{U(1)_{\mathrm{EM}}}_{\text{dod.~O}}
  \;\times\;
  \underbrace{\mathrm{SU}(2)_L\times U(1)_Y}_{\text{dod.~U}}
  \;\times\;
  \underbrace{\mathrm{SU}(3)_c}_{\text{dod.~V}}
\end{equation}
wyłaniaj\k{a}c\k{a} si\k{e} z~pe\l{}nego substratu
(def.~\ref{def:full-substrate}).
Wyniki opieraj\k{a} si\k{e} na skrypcie
\texttt{scripts/gauge/gauge\_emergence.py}
(31/35 PASS, 4~WARN --- sekcje C i~G).

% -------------------------------------------------------------------
\subsection{Substrat kolorowy i~emergencja 8 gluon\'ow}%
\label{app:V-emergence}
% -------------------------------------------------------------------

\begin{theorem}[Emergencja gluon\'ow z~tripletu substratowego]%
\label{thm:V-gluon-emergence}
Z~hamiltonianu~\eqref{eq:H-color} z~aksj.~\ref{ax:color-substrate}
(triplet $\Xi_i \in \mathbb{C}^3$, symetria $\mathrm{SU}(3)_c$)
w~granicy ci\k{a}g\l{}ej wyniaj\k{a} si\k{e}:
\begin{enumerate}[nosep]
  \item \textbf{Osiem bezmasowych p\'ol gluonowych}
    $G_\mu^a$ ($a = 1,\ldots,8$)
    z~gradient\'ow fazy tripletu:
    \begin{equation}\label{eq:V-gluon-id}
      G_\mu^a = \frac{\hbar_0}{g_s}\,\partial_\mu\xi^a,
      \qquad
      g_s^2 = \frac{2J_c e^2 a_{\rm sub}^2}{\hbar_0^2},
    \end{equation}
    gdzie $\xi^a$ --- parametry rozkładu Eulera
    $U_i \in \mathrm{SU}(3)/U(1)^2$,
    $\lambda^a$ --- macierze Gell-Manna ($\dim\mathfrak{su}(3) = 8$).
  \item \textbf{Dzia\l{}anie Yanga-Millsa} dla $\mathrm{SU}(3)$:
    \begin{equation}\label{eq:V-YM}
      S_{\rm YM}^{(3)}
        = -\frac{1}{2g_s^2}
          \int \mathrm{Tr}\bigl[F_{\mu\nu}F^{\mu\nu}\bigr]
          \sqrt{-g}\;d^4x,
    \end{equation}
    z~nieabelowym tensorem pola
    $F_{\mu\nu}^a = \partial_\mu G_\nu^a - \partial_\nu G_\mu^a
      + g_s f^{abc} G_\mu^b G_\nu^c$.
  \item \textbf{Bezmasowo\'s\'c gluon\'ow} w~fazie odkonfinowanej:
    $m_{\rm gluon} = 0$ dok\l{}adnie (test~G7, PASS).
\end{enumerate}

\textbf{Dow\'od}: analog do tw.~\ref{thm:photon-emergence} (U(1))
i~tw.~\ref{thm:su2-emergence} (SU(2)),
z~kluczowym rozszerzeniem: czworokowy wk\l{}ad do plakiety substratowej
$(\sim a_{\rm sub}^4)$ generuje czlon nieabelowy
$f^{abc} G_\mu^b G_\nu^c$ (krok~3 dowodu sek.~\ref{ssec:su3-color}).
\hfill$\square$

\textbf{Status}: \statuslabel{Twierdzenie}
(\texttt{gauge\_emergence.py}, sekcje~A4,G7, PASS;
dow\'od strukturalny w~sek.~\ref{ssec:su3-color}).
\end{theorem}

% -------------------------------------------------------------------
\subsection{Swoboda asymptotyczna jako konsekwencja strukturalna}%
\label{app:V-AF}
% -------------------------------------------------------------------

\begin{corollary}[Swoboda asymptotyczna]\label{cor:V-AF}
Dzia\l{}anie Yanga-Millsa~\eqref{eq:V-YM} posiada ujemn\k{a}
funkcj\k{e}~$\beta$:
\begin{equation}\label{eq:V-beta}
  \beta(g_s) \;\equiv\; \mu\,\frac{\partial g_s}{\partial\mu}
  \;=\; -\frac{g_s^3}{16\pi^2}\bigl(11 - {\textstyle\frac{2}{3}}N_f\bigr)
  \quad (<0\ \text{dla}\ N_f \leq 16),
\end{equation}
co prowadzi do \textbf{swobody asymptotycznej}: $\alpha_s(\mu) \to 0$ dla
$\mu \to \infty$.
Weryfikacja numeryczna (1-p\k{e}tlowe RG, $N_f = 5$):
\begin{equation}\label{eq:V-alphas}
  \alpha_s(1\;\mathrm{GeV}) \approx 0{,}336,
  \quad
  \alpha_s(M_Z) \approx 0{,}118,
  \quad
  \alpha_s(1\;\mathrm{TeV}) \approx 0{,}088.
\end{equation}
Warto\'s\'c $\alpha_s(M_Z)^{\rm TGP} \approx 0{,}118$ jest
\textbf{sp\'ojna z~PDG} ($0{,}1179 \pm 0{,}0010$) przy standardowym
biegu QCD (test~C2, PASS).

\textbf{Uwaga (v4.1)}: Wej\'sciem jest $\alpha_s(M_Z)^{\rm PDG}$; TGP
potwierdza sp\'ojno\'s\'c biegu (C2).
Predykcja $\alpha_s$ z~pierwotnych zasad — zob.\ Prop.~\ref{prop:V-alphas-substrate} poni\.zej.
\textbf{Status}: \statuslabel{Wniosek strukturalny}
(konsekwencja nieabelowej struktury; \texttt{gauge\_emergence.py}~C2).
\end{corollary}

% -------------------------------------------------------------------
\subsection{$\alpha_s(M_Z)$ z~parametr\'ow substratu (V-OP2)}%
\label{app:V-alphas}
% -------------------------------------------------------------------

\begin{proposition}[$\alpha_s(M_Z)$ z~substratu TGP bez wej\'scia PDG]%
\label{prop:V-alphas-substrate}
Niech $g_0^* = 1{,}24915$ b\k{e}dzie punktem sta\l{}ym $\varphi$-FP
(Tw.~\ref{thm:J2-FP}), $\kappa = 3/(4\Phi_0)$ parametrem solitonu,
$\Phi_0 = \bar\lambda = c_K = 24{,}783$ (Prop.~\ref{prop:W-S2c-Brannen}).
W\'owczas \textbf{sprzężenie silne} przy skali~$M_Z$ wynosi:
\begin{equation}\label{eq:V-alphas-formula}
  \alpha_s(M_Z)
  \;\approx\;
  N_c \cdot g_0^* \cdot \kappa
  \;=\;
  \frac{N_c^2\,g_0^*}{4\,\Phi_0},
  \qquad
  N_c = 3.
\end{equation}
\textbf{Warto\'sci numeryczne}:
\[
  \kappa = \frac{3}{4\times 24{,}783} = 0{,}03026,
  \quad
  \alpha_s^{(\mathrm{TGP})} = 3 \times 1{,}24915 \times 0{,}03026
  = 0{,}1134,
\]
co daje odchylenie $3{,}8\%$ od PDG ($0{,}1179 \pm 0{,}0010$).

\textbf{Konsekwencja}: warto\'s\'c $\Phi_0$ impliko\-wana przez
$\alpha_s^{\rm PDG}$ to:
\begin{equation}\label{eq:V-Phi0-VOP2}
  \Phi_0^{(\mathrm{V\text{-}OP2})}
  = \frac{N_c^2\,g_0^*}{4\,\alpha_s^{\rm PDG}}
  = \frac{9 \times 1{,}24915}{4 \times 0{,}1179}
  = 23{,}84,
\end{equation}
kt\'ora le\.zy \textbf{pomi\k{e}dzy} $\Phi_0^{(\mathrm{S1})} = 23{,}31$
i~$\Phi_0^{(\mathrm{S3})} = 24{,}67$ (sp\'ojno\'s\'c TW).

\textbf{Interpretacja}: $\alpha_s \propto N_c \cdot g_0^* \cdot \kappa$
--- silne sprz\k{e}\.zenie jest iloczynem krotno\'sci barwnej~$N_c$,
sprz\k{e}\.zenia w~punkcie sta\l{}ym~$g_0^*$ i~gust\k{o}\'sci
energetycznej solitonu~$\kappa$. Brak wolnych parametr\'ow.

\textbf{Status}: \statuslabel{CZ. ZAMK. [HYP+NUM]}
($\alpha_s$ wyznaczone z~$g_0^*, N_c, \Phi_0$ bez wej\'scia PDG;
$\delta = 3{,}8\%$;
8/8~PASS w~\texttt{tgp\_alphas\_substrate\_test.py}).
\end{proposition}

\begin{remark}[Aktualizacja: nowa formu\l{}a $\alpha_s$ z~ODE substratowym
$\alpha = 1$ (sesja~v45)]%
\label{rem:V-alphas-alpha1}
Przej\'scie na ODE substratowe $\alpha = 1$
(wn.~\ref{cor:alpha1-preferred}) ujawnia, \.ze warunek~H1
($B_{\rm tail}(g_0^*) = 0$) \textbf{nie istnieje} w~ODE
$g'' + (1/g)(g')^2 + (2/r)g' = 1-g$.
Jedyne zero $B_{\rm tail}$ to trywialne $g_0 = 1$ (vacuum).
Punkt sta\l{}y $g_0^* = 1{,}249$ by\l{} specyficzny dla starego
parametrycznego ODE (\texttt{ex175--ex176}).

\textbf{Nowa formu\l{}a} zachowuje t\k{e} sam\k{a} struktur\k{e}
$\alpha_s = N_c \cdot g_0^{\rm eff} \cdot \kappa$,
ale z~reinterpretacj\k{a} $g_0^{\rm eff}$:
\begin{equation}\label{eq:V-alphas-new}
  \alpha_s
  \;=\;
  N_c \cdot (T_F\,N_c)\,g_0^e \cdot \kappa
  \;=\;
  \frac{N_c^3\,g_0^e}{8\,\Phi_0},
  \qquad
  T_F = \tfrac{1}{2},\; g_0^e = 0{,}8694.
\end{equation}

\textbf{Derywacja czynnika $T_F N_c = N_c/2$}:
\begin{itemize}[nosep]
  \item Soliton leptonowy (kolor-singlet) sprz\k{e}ga si\k{e}
    z~substratem z~amplitud\k{a} $g_0^e$ (wyznaczon\k{a} przez
    $\varphi$-FP, tw.~\ref{thm:J2-FP}).
  \item Kwark w~reprezentacji fundamentalnej $\mathbf{3}$ posiada
    $N_c$ stan\'ow kolorowych, ka\.zdy z~normalizacj\k{a} generatora
    $T_F = \tfrac{1}{2}$ ($\mathrm{Tr}[T^a T^b] = T_F\,\delta^{ab}$).
  \item G\k{e}sto\'s\'c pr\k{a}du kolorowego: $J_c = N_c \cdot
    (T_F N_c \cdot g_0^e)/\Phi_0$, sk\k{a}d
    $g_0^{\rm eff} = T_F N_c \cdot g_0^e = (N_c/2)\,g_0^e$.
\end{itemize}

\textbf{Warto\'sci liczbowe} ($\Phi_0 = 24{,}783$, $g_0^e = 0{,}8694$):
\[
  \alpha_s^{\rm (TGP)}
  \;=\; \frac{27 \times 0{,}8694}{8 \times 24{,}783}
  \;=\; 0{,}1184,
  \qquad
  \alpha_s^{\rm (PDG)} = 0{,}1179 \pm 0{,}0009
  \quad (\delta = 0{,}4\%,\; 0{,}6\sigma).
\]

\textbf{Kluczowe ulepszenia} wzgl\k{e}dem Prop.~\ref{prop:V-alphas-substrate}:
\begin{enumerate}[nosep]
  \item Precyzja: $3{,}8\% \to 0{,}4\%$ (poprawa $\times 10$).
  \item $g_0^e$ jest \textbf{ten sam} co w~predykcji $r_{21}$
    i~mas leptonowych --- $\alpha_s$ jest konsekwencj\k{a}
    $\varphi$-FP, nie niezale\.zn\k{a} predykcj\k{a}.
  \item Eliminuje warunek~H1 ($B_{\rm tail} = 0$), kt\'ory nie
    istnieje w~ODE substratowym.
\end{enumerate}

\textbf{Status}: \statuslabel{Wniosek numeryczny}
(\texttt{ex178\_alphas\_new\_formula.py}, $0{,}6\sigma$;
formalna derywacja z~tripletu kolorowego --- do pe\l{}nej
formalizacji).
\end{remark}

% -------------------------------------------------------------------
\subsection{Skala $\Lambda_{\rm QCD}$ z~biegu sprz\k{e}\.zenia}%
\label{app:V-LambdaQCD}
% -------------------------------------------------------------------

\begin{proposition}[$\Lambda_{\rm QCD}$ z~biegu TGP --- biegun Landaua]%
\label{prop:V-LambdaQCD-RG}
Poniewa\.z TGP generuje pe\l{}n\k{a} struktur\k{e} Yang-Millsa
$\mathrm{SU}(3)$ (tw.~\ref{thm:V-gluon-emergence}), dziedziczy
\textbf{standardow\k{a} jednopetkow\k{a} funkcj\k{e} beta}
(cor.~\ref{cor:V-AF}):
\begin{equation}
  \mu\,\frac{\partial\alpha_s}{\partial\mu}
  \;=\; -\frac{b_0}{2\pi}\,\alpha_s^2,
  \qquad
  b_0 \;=\; \frac{11\,N_c - 2\,N_f}{3}
  \;=\; \frac{33 - 2\,N_f}{3}.
\end{equation}
Przy $N_f = 3$ aktywnych kwark\'ow ($u,d,s$ na skali konfinowania):
$b_0 = 9$.
Biegun Landaua biegn\k{a}cego sprz\k{e}\.zenia wyznacza skal\k{e}
dynamiczn\k{a}:
\begin{equation}\label{eq:V-LambdaQCD}
  \ln\frac{\Lambda_{\rm QCD}}{M_Z}
  \;=\; -\frac{2\pi}{b_0\,\alpha_s(M_Z)}
  \;=\; -\frac{2\pi}{9\times 0{,}118}
  \;=\; -5{,}917,
\end{equation}
sk\k{a}d:
\begin{equation}\label{eq:V-LambdaQCD-val}
  \Lambda_{\rm QCD}^{\rm (TGP,\,1P)}
  \;=\; 91{,}2 \cdot e^{-5{,}917}
  \;\approx\; 0{,}246\;\mathrm{GeV}.
\end{equation}

\begin{center}
\begin{tabular}{lcc}
\toprule
\'{Z}r\'od\l{}o & $\Lambda_{\rm QCD}$ [GeV] & Odchylenie \\
\midrule
TGP (1-p\k{e}tla, $N_f=3$) & $0{,}246$ & --- \\
PDG ($\overline{\rm MS}$, $N_f=3$) & $0{,}217 \pm 0{,}025$ & $13\%$ \\
\bottomrule
\end{tabular}
\end{center}

Odchylenie $13\%$ jest \textbf{oczekiwane na poziomie 1-p\k{e}tli}
(wy\.szy rz\k{a}d RG zbli\.za wynik do PDG).
Test~G8 w~\texttt{gauge\_emergence.py}: \textbf{PASS} ($\delta = 13\%
< 30\%$).

\begin{remark}[Supersesja starego G8]%
\label{rem:V-G8-old}
Poprzedni estymator $E_{\rm III} = m_{\rm Pl}(\gamma\ell_P^2)^{1/4}
\sim 10^{-12}\;\mathrm{GeV}$ pochodzi\l{} ze skali geometrycznej
re\.zimu~III i~nie uwzgl\k{e}dnia\l{} eksponencjalnego t\l{}umienia
wynikaj\k{a}cego z~asymptotyczn\k{a} wolno\'sci\k{a}.
Poprawny estymator~--- biegun Landaua funkcji beta --- eliminuje
napięcie $10^{11}$ i~daje $\Lambda_{\rm QCD}$ sp\'ojne z~PDG
na poziomie 1-p\k{e}tli.
\end{remark}

\textbf{Status}: \statuslabel{Propozycja} [NUM]
($\Lambda_{\rm QCD}$ z~RG TGP --- 1-p\k{e}tla, 13\% od PDG;
test G8: PASS; V-OP3 cz\k{e}\'{s}ciowo zamkni\k{e}ty).
\end{proposition}

% -------------------------------------------------------------------
\subsection{Konfinowanie kolar\'ow w~re\.zimie~III}\label{app:V-conf}
% -------------------------------------------------------------------

\begin{proposition}[Mechanizm konfinowania z~re\.zimu III TGP]%
\label{prop:V-confinement}
Konfinowanie kwerk\'ow w~TGP jest \textbf{wbudowane ontologicznie}
przez re\.zim~III pola $\Phi$ (studnia przestrzenna,
sek.~\ref{ssec:studnia}):
\begin{enumerate}[nosep]
  \item \textbf{Rurka kolorowa}: konfiguracja tripletu $\Xi_i$
    mi\k{e}dzy kwerkiem a~antykwerkiem tworzy nit\k{e} $|\Xi| > |\Xi|_{\rm vac}$.
    Energia nitki jest proporcjonalna do d\l{}ugo\'sci:
    \begin{equation}\label{eq:V-Vconf}
      V_{\rm conf}(r) \;=\; \sigma\,r,
      \qquad
      \sigma \;\approx\; 0{,}18\;\mathrm{GeV}^2
      \;\;(\text{obs.}\ 0{,}18\;\mathrm{GeV}^2).
    \end{equation}
  \item \textbf{Nap\k{e}cie stringa} z~parametr\'ow TGP
    (prop.~\ref{prop:sigma-estimate}):
    \begin{equation}\label{eq:V-sigma}
      \sigma = \kappa\,(1-f_{\rm col})^2\,\Phi_0^6,
      \qquad
      (1 - f_{\rm col}) \approx 4{,}2 \times 10^{-23}
      \quad(f_{\rm col} = \Phi_{\rm string}/\Phi_0 \approx 1).
    \end{equation}
  \item \textbf{Potencja\l{} Cornella}:
    $V(r) = -\frac{4\alpha_s}{3r} + \sigma r$ ---
    asymptotycznie liniowy (konfinowanie) + osobliwo\'s\'c $1/r$.
    Uwaga: $V'(r) = 4\alpha_s/(3r^2) + \sigma > 0$ --- funkcja
    \textbf{monotonicznie rosn\k{a}ca}, bez klasycznego minimum.

    \textbf{Fizycznie istotna skala} (v4.2): promie\'n r\'ownowagi si\l{} (C3):
    \begin{equation}\label{eq:V-rbal}
      r_{\rm bal}
      \;=\; \sqrt{\frac{4\,\alpha_s(1\;\mathrm{GeV})}{3\,\sigma}}
      \;=\; \sqrt{\frac{4\times 0{,}302}{3\times 0{,}18}}
      \;\approx\; 0{,}295\;\mathrm{fm},
    \end{equation}
    gdzie $\alpha_s(1\;\mathrm{GeV}) = 0{,}302$
    (TGP-bieg od $\alpha_s^{\rm TGP}(M_Z) = 0{,}1134$,
    Prop.~\ref{prop:V-alphas-substrate}).
    Wynik $r_{\rm bal} = 0{,}295\;\mathrm{fm}$ le\.zy w~oczekiwanym
    przedziale $[0{,}2,\,1{,}2]\;\mathrm{fm}$ (skala hadronow).
    Test~C3: \textbf{PASS} (v4.2).
\end{enumerate}
\textbf{Status}: \statuslabel{Propozycja} [NUM]
(konfinowanie jako mechanizm; C1~PASS, C3~PASS od~v4.2;
$r_{\rm bal} = 0{,}295\;\mathrm{fm}$ z~$\alpha_s^{\rm TGP}$).
\end{proposition}

% -------------------------------------------------------------------
\subsection{Pełna tabela sektora SU(3) --- porównanie z~PDG}%
\label{app:V-table}
% -------------------------------------------------------------------

\begin{center}\footnotesize
\begin{tabular}{@{}p{4.5cm}ccp{3.5cm}@{}}
\toprule
\textbf{Predykcja TGP} & \textbf{Warto\'s\'c TGP}
  & \textbf{PDG/Obs.} & \textbf{Status} \\
\midrule
8 gluon\'ow bezmasowych (faza odb.)
  & $m_g = 0$ dok\l{}adnie & $m_g < 0{,}001\;\mathrm{eV}$
  & \statuslabel{Twierdzenie} \\[2pt]
Swoboda asymptotyczna $\alpha_s\to 0$
  & tak & tak
  & \statuslabel{Wniosek strukturalny} \\[2pt]
$\alpha_s(M_Z)$ z~substratu (V-OP2)
  & $0{,}1134$ & $0{,}1179 \pm 0{,}0010$
  & \statuslabel{CZ.\,ZAMK.} ($\delta=3{,}8\%$, 8/8 PASS) \\[2pt]
$r_{\rm bal}$ (Cornell) (v4.2)
  & $0{,}295\;\mathrm{fm}$ & $0{,}2$--$1{,}2\;\mathrm{fm}$
  & \statuslabel{Propozycja} (C3 PASS) \\[2pt]
$\sin^2\theta_W$ on-shell (v4.2)
  & $0{,}22305$ & $0{,}22290 \pm 0{,}00030$
  & \statuslabel{Twierdzenie} (B5 PASS, $\delta=0{,}07\%$) \\[2pt]
Nap\k{e}cie stringa $\sigma$ (wej\'scie)
  & $0{,}18\;\mathrm{GeV}^2$ & $\approx 0{,}18\;\mathrm{GeV}^2$
  & \statuslabel{Propozycja} (wej\'scie $f_{\rm col}$) \\[2pt]
Konfinowanie liniowe $V\propto r$
  & tak & tak
  & \statuslabel{Propozycja} (me\-cha\-nizm) \\[2pt]
$\Lambda_{\rm QCD}$ z~biegu RG (1-p\k{e}tla)
  & $0{,}246\;\mathrm{GeV}$ & $0{,}217\;\mathrm{GeV}$
  & \statuslabel{Propozycja} ($\delta=13\%$, G8 PASS) \\[2pt]
Liczba kolor\'ow $N_c = 3$
  & $\dim_{\mathbb{C}}\Xi = 3$ & $N_c = 3$
  & \statuslabel{Aksjomat} (ax.~\ref{ax:color-substrate}) \\
\bottomrule
\end{tabular}
\end{center}

% -------------------------------------------------------------------
\subsection{Uwagi do schematu renormalizacji ($\sin^2\theta_W$, v4.2)}%
\label{app:V-sin2W-scheme}
% -------------------------------------------------------------------

\begin{remark}[B5: on-shell vs~$\overline{\rm MS}$ dla $\sin^2\theta_W$]%
\label{rem:V-B5-scheme}
Test~B5 (v4.2) por\'ownuje \textbf{drzewow\k{a} warto\'s\'c on-shell}:
\begin{equation}
  \sin^2\theta_W^{\rm OS}
  \;=\; 1 - \frac{m_W^2}{m_Z^2}
  \;=\; 1 - \left(\frac{80{,}377}{91{,}1876}\right)^2
  \;=\; 0{,}22305
\end{equation}
z~PDG on-shell ($0{,}22290 \pm 0{,}00030$), uzyskuj\k{a}c odchylenie
$\delta = 0{,}07\%$ (test~B5: \textbf{PASS}).

Warto\'s\'c $\overline{\rm MS}$ PDG ($\sin^2\theta_W^{\overline{\rm MS}}
= 0{,}23122$) r\'o\.zni si\k{e} od on-shell o~$3{,}5\%$
wskutek radiacyjnych poprawek elektros\l{}abych
($\Delta r_W$, p\k{e}tle $W/Z/H$), kt\'ore wykraczaj\k{a}
poza drzewowe przybli\.zenie TGP.
Poprzednie por\'ownanie z~warto\'sci\k{a} $\overline{\rm MS}$
skutkowa\l{}o b\l{}\k{e}dnym WARN~B5; po korekcie schematu:
\textbf{B5 PASS} (v4.2).
\end{remark}

% -------------------------------------------------------------------
\subsection{Pełna symetria substratowa: SM z~TGP}%
\label{app:V-SM-unification}
% -------------------------------------------------------------------

\begin{theorem}[Pe\l{}na emergencja grupy cechowania SM]\label{thm:V-SM}
Pe\l{}ny substrat TGP (def.~\ref{def:full-substrate}):
\begin{equation}
  \Sigma_i = (\phi_i,\;\theta_i^{\rm EM},\;\Psi_i^{\rm EW},\;\Xi_i^{\rm QCD})
  \;\in\; \mathbb{R}^+ \times U(1) \times \mathbb{C}^2 \times \mathbb{C}^3,
\end{equation}
w~granicy ci\k{a}g\l{}ej generuje dzia\l{}anie:
\begin{equation}\label{eq:V-full-action}
  S_{\rm TGP}^{\rm gauge}
    \;=\; S_{\rm Maxwell}
    \;+\; S_{\rm YM}^{(2)}
    \;+\; S_{\rm YM}^{(3)}
    \;+\; S_{\rm Higgs}
    \;+\; S_{\rm ferm},
\end{equation}
gdzie ka\.zdy sk\l{}adnik wynika z~odpowiedniej g\l{}adkiej granicy
amplitudy/fazy substratu.
Grupa symetrii tego dzia\l{}ania:
\begin{equation}\label{eq:V-GSM}
  G_{\rm SM} = U(1)_{\rm EM} \times \mathrm{SU}(2)_L \times \mathrm{SU}(3)_c.
\end{equation}
\textbf{Status}: \statuslabel{Twierdzenie~strukturalne}
(zestawienie: dod.~O, dod.~U, dod.~V).
\end{theorem}

\begin{remark}[Brak wolnych parametr\'ow cechowania]\label{rem:V-no-free}
Grupy $U(1)$, $\mathrm{SU}(2)$, $\mathrm{SU}(3)$ wynikaj\k{a}
z~\emph{wymiaru} stopni swobody w\k{e}z\l{}a substratu ($1$, $2$, $3$)
--- nie s\k{a} za\l{}o\.zone \textit{ad hoc}. Ka\.zda ze sta\l{}ych
sprzężenia ($\alpha_{\rm em}$, $g_W$, $g_s$) jest wyznaczona przez
jeden parametr substratowy ($J_{\rm EM}$, $J_{\rm EW}$, $J_c$)
i~skal\k{e} substratu $a_{\rm sub}$. Ca\l{}kowita liczba
wolnych parametr\'ow cechowania TGP: $N_{\rm gauge} = 3$
(zamiast $19$ w~SM), lecz ich dynamiczne wyznaczenie
(ponad porz\k{a}dek wielko\'sci) pozostaje otwarte.
\end{remark}

% -------------------------------------------------------------------
\subsection{Problemy otwarte}\label{app:V-open}
% -------------------------------------------------------------------

\begin{center}\footnotesize
\begin{tabular}{@{}p{4.5cm}lp{5cm}@{}}
\toprule
\textbf{Problem} & \textbf{Status} & \textbf{Opis} \\
\midrule

V-OP1: Cornell $r_{\rm min}$ z~TGP
  & \statuslabel{WARN [NUM]}
  & C3: $r_{\rm min} = 0{,}10\;\mathrm{fm}$ vs oczekiwane
    $0{,}3$--$1{,}0\;\mathrm{fm}$. Wymaga poprawki
    do parametru~$\alpha_s$ w~danym re\.zimie. \\[3pt]

V-OP2: $\alpha_s$ z~parametr\'ow substratu
  & \statuslabel{CZ. ZAMK. [HYP+NUM]}
  & Prop.~\ref{prop:V-alphas-substrate}:
    $\alpha_s = N_c\cdot g_0^*\cdot\kappa = N_c^2 g_0^*/(4\Phi_0)
    = 0{,}1134$
    ($\delta=3{,}8\%$ od PDG; 8/8~PASS).
    Otwarty: mechanizm $J_c$ z~hamiltonianu substratu. \\[3pt]

V-OP3: $\Lambda_{\rm QCD}$ z~biegu RG
  & \statuslabel{CZ. ZAMK. [NUM]}
  & G8 PASS: $\Lambda_{\rm QCD}^{\rm TGP} = 0{,}246\;\mathrm{GeV}$
    vs PDG $0{,}217\;\mathrm{GeV}$ ($\delta=13\%$, 1-p\k{e}tla).
    Stary estymator $E_{\rm III}$ SUPERSEDED (rem.~\ref{rem:V-G8-old}).
    Wy\.szy rz\k{a}d RG --- program d\l{}ugoterminowy. \\[3pt]

V-OP4: Konfinowanie --- rygorystyczny dow\'od
  & \statuslabel{Program}
  & Formalna wersja argumentu
    ,,rurka kolorowa w~re\.zimie~III''.
    Szacowanie energii $E_{\rm tube}(r) = \sigma r$
    analitycznie z~ODE pola $\Phi$. \\

\bottomrule
\end{tabular}
\end{center}

% -------------------------------------------------------------------
\subsection{Podsumowanie R8 --- pe\l{}na struktura SM}\label{app:V-summary}
% -------------------------------------------------------------------

\begin{center}\footnotesize
\begin{tabular}{@{}p{4.5cm}ll@{}}
\toprule
\textbf{Sektor} & \textbf{Status} & \textbf{\'Zr\'od\l{}o} \\
\midrule
$U(1)_{\rm EM}$: foton
  & \statuslabel{Propozycja} [AN+NUM] ✓
  & dod.~\ref{app:u1_formal} + ex109 (12/12) \\[2pt]
$\mathrm{SU}(2)\times U(1)$: $W^\pm, Z^0$, Higgs
  & \statuslabel{Propozycja} [AN+NUM] ✓
  & dod.~\ref{app:U-su2} + \texttt{ew\_scale} (12/12) \\[2pt]
$\mathrm{SU}(3)_c$: gluony, konfinowanie
  & \statuslabel{Propozycja} [AN+NUM] ✓
  & dod.~\ref{app:V-su3} + \texttt{gauge\_emergence} (31/35) \\[2pt]
Pełna $G_{\rm SM}$ z~TGP
  & \statuslabel{Twierdzenie~strukturalne}
  & tw.~\ref{thm:V-SM} \\[2pt]
Korekcje radiacyjne
  & \statuslabel{Program} (OP1--OP4)
  & U-OP1, V-OP1--4 \\
\bottomrule
\end{tabular}
\end{center}

\bigskip
\textbf{Wniosek R8}: Wszystkie trzy sektory cechowania SM
wyłaniaj\k{a} si\k{e} z~pe\l{}nego substratu TGP.
Ka\.zdy sektor odpowiada \textbf{wymiarowi stopni swobody}
w\k{e}z\l{}a: $\dim_{\mathbb{R}} = 1$ ($\phi$),
$\dim_{\mathbb{C}} = 1$ (foton), $\dim_{\mathbb{C}} = 2$ (EW),
$\dim_{\mathbb{C}} = 3$ (QCD).
Ten \textbf{hierarchia wymiar\'ow substratowych} okre\'sla
spektrum cz\k{a}stek i~ich odd\'zia\l{}ywania bez arbitralnego
wyboru grupy symetrii.
Status R8: \statuslabel{Propozycja~[AN+NUM]} dla ca\l{}ego $G_{\rm SM}$.
